\section{The imported theorem stacks}
\label{sec:imports}

The closure separates its inputs into two classes.  The present section
records the imported theorem stacks: established arithmetic geometry
that the closure consumes as proof-carrying hypotheses.  The next section
records the recognition sources: arithmetic content carried by the
formed BSD layer.  This distinction matters.  The imported stacks are not
new theorems here, and the closure never derives them
from the framework.  It derives consequences from them, after their
provenance and downstream role have been made explicit.

There are three headline imports in this section.  The first controls the
Cassels--Tate side of the finite \(p\)-primary comparison.  The second
supplies the Selmer-complex pairing substrate \( \AofE \).  The third
supplies the rank-\(\geq 2\) archimedean determinant comparison.  Together
they keep the classical arithmetic geometry used by the landing theorem out
of the framework argument while still making clear exactly which classical
inputs are being used.  The overconvergent \(\eta\) and
\(T_{\mathrm{CASCADE}}\) imports are introduced later as supporting
conditional stacks and are counted separately in App~\ref{app:formalization}.

The symbol \(\chiCTp\) refers here to a Cassels--Tate comparison: a
comparison between a square-root Fitting object on a Selmer-complex
cohomology group and the Pfaffian scalar attached to the Cassels--Tate
pairing.  In ordinary terms, this is the input that lets the \(p\)-primary
finite piece be read with the correct orientation and sign.

\begin{definition}[Imported $\chi_{CT,p}$ theorem stack]\label{def:closure:chi-ct-p-import}
A proof-carrying structure bundling the eight named external inputs
\(T_{E1}\ldots T_{E8}\): Sakamoto's Stark-system theory~\citep{Sakamoto2018};
the determinant/Stark-system and derived-height comparisons of Macias
Castillo--Sano~\citep[Thms.~3.4 and~4.2]{MaciasCastilloSano2026};
Knudsen--Mumford determinant theory~\citep{KnudsenMumford1976};
Nekov\'a\v{r}'s Selmer-complex pairing~\citep[\S10.8.7]{Nekovar2006};
Flach's pairing comparison~\citep{Flach1990}; and the classical
2-Selmer normalization of Fisher--Schaefer--Stoll~\citep[Thm.~6.3]{FisherSchaeferStoll2010}.
The results are stored as proposition fields,
plus the unconditional sign closure
\(\Sigma_{\mathrm{NekCT}}=+1\).  It is taken as a hypothesis by
\Cref{thm:closure:chi-ct-p-comparison}; it is not derived in this
paper.\footnote{Encoded in Lean as a typed structure carrier
\texttt{SixBirdsBSD.\allowbreak Closure.\allowbreak Imports.\allowbreak chiCTpImport}
bundling the imported theorems as hypothesis fields; not a Lean axiom.
See App~\ref{app:formalization}.}
\end{definition}

The downstream role of this stack is narrow.  It is not a general
replacement for the Cassels--Tate literature.  It is the specific
orientation-respecting package needed for the \(\chiCTp\) theorem:
given the usual finite \(\Sha[p^\infty]\), nondegeneracy, and
Selmer-complex hypotheses, it supplies the Pfaffian comparison and the
sign closure \(\Sigma_{\mathrm{NekCT}}=+1\).

The second import packages the Selmer-complex pairing substrate denoted
\(\AofE\).  This notation is local to this article: \(\AofE\) is the
unified Selmer-complex pairing object, not the analytic leading coefficient
written \(A_E\) in \citet{TsiokosBSDApparatus2026}.  Mathematically, it is the
ambient category, Selmer complex, and pairing target in which the
Cassels--Tate pairing is intrinsic and compatible with the
determinant-line formulation.

\begin{definition}[Imported $\AofE$ Selmer-complex pairing]\label{def:closure:a-e-import}
A proof-carrying structure for the \(\AofE\) unified Selmer-complex
substrate (Nekov\'a\v{r} Selmer complex~\citep{Nekovar2006} with intrinsic Cassels--Tate
pairing; Burns--Flach and related Macias Castillo--Sano--Nekov\'a\v{r}
ETNC-adjacent input~\citep{BurnsFlach2001,BlochKato1990,FontainePerrinRiou1994,BurnsMaciasCastilloSeo2024,MaciasCastilloSano2026}
\(\AofE^{\mathrm{Sel}}\)): a typed target category with a
bilinear/categorical pairing field.  It is used only as a
hypothesis.\footnote{Encoded in Lean as a typed structure carrier
\texttt{SixBirdsBSD.\allowbreak Closure.\allowbreak Imports.\allowbreak aEImport}
bundling the imported theorem as a hypothesis field; not a Lean axiom.
See App~\ref{app:formalization}.}
\end{definition}

This import is the bridge between two ways of seeing the same finite
arithmetic information.  On one side is the Selmer-complex description
with its intrinsic Cassels--Tate pairing.  On the other is the
determinant-line language used in the closure landing.  The \(\AofE\)
carrier records that the target category and pairing data have already
been supplied, so later arguments may use the pairing without
reconstructing the Selmer-complex theory.

The third import is archimedean.  In rank at least two, the landing
needs a comparison between the Beilinson determinant side and the
leading \(L\)-value side.  This is where the regulator-to-determinant
identification enters.  The closure treats it as a classical imported
comparison, separate from the higher-rank recognition source that will
encode the remaining Gross--Zagier-type fixity.

\begin{definition}[Imported Beilinson rank-$\geq 2$ comparison]\label{def:closure:beilinson-import}
A proof-carrying structure for the Beilinson determinant/\(L\)-value~\citep{Beilinson1985}
comparison in rank \(r\geq 2\): the archimedean
regulator-to-determinant-line identification used by the rank-\(\geq 2\)
landing.  It is used only as a hypothesis.\footnote{Encoded in Lean as a
typed structure carrier
\texttt{SixBirdsBSD.\allowbreak Closure.\allowbreak Imports.\allowbreak beilinsonImport}
bundling the imported theorem as a hypothesis field; not a Lean axiom.
See App~\ref{app:formalization}.}
\end{definition}

The point of isolating this comparison is to prevent the closure
argument from hiding an archimedean theorem inside the recognition
predicate.  The Beilinson import supplies the determinant/\(L\)-value
identification; the recognition source introduced later supplies the
formed-layer fixity condition that makes the rank-\(\geq 2\) readout
land on the Strong-BSD scalar.

\begin{nonclaim}\label{nonclaim:closure:imports-not-derived}
The \(\chiCTp\) comparison remains conditional on its imported stack
\(T_{E1}\ldots T_{E8}\); the import structures are sourced assumptions,
not results of this paper.  The closure derives consequences
\emph{from} them and never proves them.  They are carried as typed
hypothesis structures and introduce no project-local axioms.
\end{nonclaim}

% ===========================================================================
